\section{Introduction}

Systems that translate policies, contracts, or statutes into executable logic
create two superficially similar artifacts: an interpretation of a source and a
program that can be run. Passing a type checker, solver, or execution test says
something about the second artifact, but not necessarily the first. A valid
decision table can omit an exception; two engines can agree on the same
mistranslation; and high conditional accuracy can be obtained by refusing the
hardest cases without reporting coverage.

Existing resources measure important pieces of this problem. LegalBench and
RuleArena test legal and rule-guided reasoning
\citep{guha2023legalbench,zhou2025rulearena}; ContractNLI and CUAD pair legal
text with labels or evidence spans \citep{koreeda2021contractnli,hendrycks2021cuad};
and recent work evaluates generated decision models by structural and outcome
equivalence \citep{graus2026legaltext}. Yet a system-level evaluation can
still conflate source support, executable behavior, and pipeline operation, or
silently omit refused and unreproducible cases.

We ask a narrower question: \emph{what evidence is required before an
executable artifact can support a claim about its source?} Our answer is an
evidence-gated protocol in which every quantity names its observation unit,
denominator, independent anchor, and epistemic state. The protocol treats
refusal as an observable outcome and reports executable yield separately from
quality conditional on execution. Failed replays and unavailable labels stay
in the evidence ledger rather than becoming missing rows.

We instantiate the protocol in \lexec{} (Lexically grounded EXecutable
Compliance), a bounded text-to-logic compiler and review instrument. Its typed
\lexir{} makes the supported semantic subset explicit; lowering is total and
fail-closed; and provenance, validation, execution, and reviewer records remain
linked by hashes. The compiler is not presented as a legal-reasoning oracle. It
is the apparatus with which source fidelity and executable behavior can be
measured without confusing one for the other.

This paper makes three contributions:
\begin{enumerate}
  \item \textbf{An evidence-gated evaluation contract.} We formalize three
  evidence planes---attributable source fidelity (AFS), operational evidence
  (OPS), and outcome equivalence (OE)---and define claim-promotion rules that
  retain refusals, invalid runs, and missing anchors.
  \item \textbf{A provenance-preserving compiler instrument.} We describe an
  implemented bounded IR, deterministic validators, FEEL/reference evaluator,
  DMN emitter, solver diagnostics, run manifests, and review workbench.
  \item \textbf{A reproducibility audit with visible negative evidence.} In a
  retained 1,900-row benchmark replay, 792 rows (41.7\%) disagree with the
  released aggregate. We quarantine the anchor and report the disagreement
  structure instead of selecting the favorable score. A separate
  1,012-document run demonstrates operational coverage while remaining
  explicitly non-claiming about semantic quality.
\end{enumerate}

The result is an evaluation artifact and a concrete audit finding, not a claim
of population-level legal correctness. The latter requires licensed source
annotations, independent execution, and a frozen human-evaluation protocol;
those gates are stated explicitly and are not replaced by unit tests or
exploratory throughput.
